\newcommand\expror[2]{\ensuremath{\texttt{Either}\ #1\ #2}\xspace}
\newcommand\exprand[2]{\ensuremath{(#1, #2)}\xspace}
\newcommand\exprorleft[1]{\ensuremath{\texttt{Left}\ #1}\xspace}
\newcommand\exprorright[1]{\ensuremath{\texttt{Right}\ #1}\xspace}
\newcommand\tyand[2]{\ensuremath{(#1, #2)}\xspace}
\newcommand\tyor[2]{\ensuremath{\texttt{Either}\ #1\ #2}\xspace}
\newcommand\tyex[3]{\ensuremath{(#1::#2,#3)}\xspace}
\newcommand\tyall[3]{\ensuremath{#1:#2 \rightarrow #3}\xspace}

\section{Embedding Natural Deduction with Refinement Types}
\label{sec:higher-order}

In this section we show how user-provided
\emph{quantified} specifications
can be naturally encoded using
$\lambda$-abstractions and
dependent pairs to encode
universal and existential
quantification, respectively.
Proof terms can be generated
using the standard natural
deduction derivation rules,
following Propositions as Types~\cite{Wadler15}
(also known as the Curry-Howard
isomorphism~\cite{CHIso}).
%
What is new is that we exploit this encoding to show for the first time that a refinement
type system can represent any proof
in Gentzen's natural deduction~\cite{Gentzen35}
while still taking advantage of
SMT decision procedures to automate the
quantifier-free portion of natural
deduction proofs.
%
For simplicity, in this section we assume
all terms are total.

\subsection{Propositions: Refinement Types}

Figure~\ref{fig:hol} maps logical predicates
to types constructed over quantifier-free refinements.

\mypara{Native terms}
%
Native terms consist of all of the
(quantifier-free) expressions
of the refinement languages.
%
In \S~\ref{sec:formalism}
we formalize refinement typing
in a core calculus \corelan where
refinements include (quantifier-free)
terminating expressions.


\mypara{Boolean connectives}
Implication $\formula_1 \Rightarrow \formula_2$
is encoded as a \emph{function} from the proof of
$\formula_1$ to the proof of $\formula_2$.
%
Negation is encoded as an implication where
the consequent is "False".
%
Conjunction $\formula_1 \land \formula_2$
is encoded as the \emph{pair} "($\formula_1$, $\formula_2$)"
that contains the proofs of \emph{both} conjuncts
%
and disjunction $\formula_1 \lor \formula_2$
is encoded as the \emph{sum} type "Either"
that contains the proofs of \emph{one of}
the disjuncts, \ie where "data Either a b = Left a | Right b".

\mypara{Quantifiers}
%
Universal quantification $\forall x. \formula$
is encoded as lambda abstraction $x:\typ \rightarrow \formula$
and eliminated by function application.
%
Existential quantification $\exists x. \formula$
is encoded as a dependent pair "(x::$\typ$, $\formula$)"
that contains the term $x$ and a proof of a formula
that depends on $x$.
%
Even though refinement type systems do not
traditionally come with explicit syntax for
dependent pairs, one can encode dependent
pairs in refinements using abstract
refinement types~\cite{Vazou13}
which do not add extra complexity
to the system.
%
Consequently, we add the syntax for
dependent pairs in Figure~\ref{fig:hol}
as syntactic sugar for abstract refinements.

\begin{figure}[t!]
$$
\begin{array}{rcccc}
                    & \quad  & \textbf{Logical Formula} & \quad & \textbf{Refinement Type}  \\
\text{Native Terms} & \quad&   $e$      & & \proofterm{e}    \\
\text{Implication}  & & \formula_1 \Rightarrow \formula_2 &      & {\formula_1} \rightarrow {\formula_2}     \\
\text{Negation}     & & \neg \formula                  & & \formula \rightarrow \proofterm{\efalse} \\
\text{Conjunction}  & & \formula_1 \land \formula_2    & & \tyand{\formula_1}{\formula_2} \\
\text{Disjunction}  & &  \formula_1 \lor \formula_2    & & \tyor{\formula_1}{\formula_2}   \\
\text{Forall}       & &\forall x. \formula             & & \tyall{x}{\typ}{\formula}                \\
\text{Exists}       & & \exists x. \formula            & & \tyex{x}{\typ}{\formula}                 \\
\end{array}
$$
\caption{Mapping from logical predicates to quantifier-free refinement types.
\proofterm{e} abbreviates \prooftermlarge{e}. Function binders are not relevant for
negation and implication, and hence, elided.}
\figrule
\label{fig:hol}
\end{figure}


\subsection{Proofs: Natural Deduction}

We overload \formula to be both a proposition and
a refinement type.
%
We connect these two meanings of \formula by
using the Propositions as Types~\cite{Wadler15},
to prove that if there exists an expression (or proof term) with
refinement type \formula, then the proposition
\formula is valid.
%

We construct proofs terms using Gentzen's
natural deduction system~\cite{Gentzen35},
whose rules map directly to refinement type
derivations.
%
The rules for natural deduction arise from the propositions-as-types reading
of the standard refinement type
checking rule (to be defined in
\S~\ref{sec:formalism}) \ch{\env}{e}{\formula}
as ``\formula is provable under the assumptions of \env''.
%% and
%% \ch{\env}{e}{\typ} as ``$e$ is a term''.
%
We write \gentzenValid{\env}{\formula}
for Gentzen's natural deduction judgment
``under assumption \env, proposition \formula holds''.
%
Then, each of Gentzen's logical rules
can be recovered from the rules in
Figure~\ref{fig:ch} by rewriting
each judgment \ch{\env}{e}{\formula}
of \corelan as \gentzenValid{\env}{\formula}.
%
For example, conjunction and universal
elimination can be derived as:
%
$$
\inference{
	\gentzenValid{\env}{{\formula_1} \lor{\formula_2}} &&
	\gentzenValid{\env,\formula_1}{\formula} &&
	\gentzenValid{\env,\formula_2}{\formula} &&
}{
	\gentzenValid{\env}{\formula}
}
\lor\small{\textsc{\!-E}}
\quad
\inference{
	\gentzenValid{\env}{e_x\ \text{term}} &&
	\gentzenValid{\env}{\forall x. \formula}
}{
	\gentzenValid{\env}{\formula\subst{x}{e_x}}
}
\forall\small{\textsc{\!-E}}
$$

\mypara{Programs as Proofs}
%
As Figure~\ref{fig:ch} directly maps
natural deduction rules to derivations
that are accepted by refinement typing,
we conclude that if there exists a
natural deduction derivation for a
proposition \formula, then there
exists an expression that has the refinement type \formula.
\begin{theorem}\label{theorem:embedding}
If\ \gentzenValid{\env}{\formula},
then we can construct an $e$ such that\ \ch{\env}{e}{\formula}.
\end{theorem}

Note that our embedding is \textit{not} an isomorphism, 
since the converse of Theorem~\ref{theorem:embedding} does not hold. 
%
As a counterexample, the law of the excluded middle
can be proved in our system
(\ie we can construct an "Either" term $e$, so that 
$\ch{\bind{p}{\proofterm{\typbool\ | \ \etrue}}}{e}{p\lor\neg p}$), 
but cannot be proved using natural deduction
(\ie $\gentzenValid{{\proofterm{\etrue}}\not}{p\lor\neg p}$).
%
The reason for that is that our system is using the classical logic
of the SMTs, which includes the law of the excluded middle.
On the contrary, in intuitionistic systems that also encode natural deduction
(\eg \coq, \idris, \nuPRL)
the law of the excluded middle should be axiomatized.

%
\subsection{Examples}\label{sec:expressiveness:examples}
%
Next, we illustrate our encoding
with examples of proofs for
quantified propositions
ranging from
%
textbook logical tautologies,
%
properties of datatypes like lists,
%
and induction on natural numbers.
%
%These and more examples can be found in \url{https://github.com/nikivazou/LiquidHOL}.

\mypara{Natural Deduction as Type Derivation}
%
We illustrate the mapping from natural deduction rules to typing rules
in Figure~\ref{fig:ex:ch} which uses typing judgments to express
Gentzen's proof of the proposition
%
$$
\phi \equiv (\exists x. \forall y. (p\ x\ y)) \Rightarrow (\forall y. \exists x. (p\ x\ y))
$$
%
Read bottom-up, the derivation provides a proof of \formula.
Read top-down, it constructs a \emph{proof} of the formula
as the \emph{term}
%the proof term of the formula to be
$\lambda e\ y. \texttt{case}\ e\ \texttt{of}\ \{
    (x,e_x) \rightarrow (x, e_x\ y)\}$.
%
This proof term corresponds directly to the following Haskell expression
that typechecks with type \formula.
\begin{mcode}
    exAll :: p:(a -> a -> Bool) -> (x::a, y:a -> {p x y}) -> y:a -> (x::a, {p x y})
    exAll _ = \e y -> case e of {(x, ex) -> (x, ex y)}
\end{mcode}
%
% Pattern matching simplifies the proof to: "existsForall f (x,px) y = (x, px y)".

\begin{figure}
%\begin{framed}
$$
\inference{
	\inference{
		\inference{
			\bind{e}{\formulap},
			\bind{y}{\typ_y}
			\vdash
			 e
			 :
	        \formulap &&
	        \inference{
	        \bind{e}{\formulap},
			\bind{y}{\typ_y},
			\bind{x}{\type_x},
			\bind{e_x}{\formulax}
			\vdash
			 e_x
			 :
	         \formulax
	         \\
	        \bind{e}{\formulap},
			\bind{y}{\typ_y},
			\bind{x}{\type_x},
			\bind{e_x}{\formulax}
			\vdash
			 y
			 :
	         \typ_y
	        }{
	        \bind{e}{\formulap},
			\bind{y}{\typ_y},
			\bind{x}{\type_x},
			\bind{e_x}{\formulax}
			\vdash
			 e_x\ y
			 :
	         p\ x\ y
	        }\forallElim
		}{
			\bind{e}{\formulap},
			\bind{y}{\typ_y}
			\vdash
			\texttt{case}\ e\ \texttt{of}\ \{
                                (x,e_x) \rightarrow (x, e_x\ y)
                        \}
                        :
	        \exists x. (p\ x\ y)
		}\existsElim
	}{
	\bind{e}{\formulap}
	\vdash
	\lambda y. \texttt{case}\ e\ \texttt{of}\ \{
    (x,e_x) \rightarrow (x, e_x\ y)
    \}
    :
	\forall y. \exists x. (p\ x\ y)
	}\forallIntro
}
{
\emptyset
\vdash
\lambda e\ y. \texttt{case}\ e\ \texttt{of}\ \{
    (x,e_x) \rightarrow (x, e_x\ y)
\}
:
{(\exists x. \forall y. (p\ x\ y)) \Rightarrow (\forall y. \exists x. (p\ x\ y))}
}
\impIntro
$$
\caption{Proof of $(\exists x. \forall y. (p\ x\ y)) \Rightarrow (\forall y. \exists x. (p\ x\ y))$
where
$\formulap \equiv \exists x. \forall y. (p\ x\ y)
$,
$
\formulax \equiv \forall y. (p\ x\ y)
$.
}
\figrule
\label{fig:ex:ch}
% \end{framed}
\end{figure}

\mypara{SMT-aided proofs}
The great benefit of using refinement types
to encode natural deduction is that the quantifier-free
portions of the proof can be automated via SMTs.
%
For every quantifier-free proposition \formula,
you can convert between $\{\formula\}$,
where \formula is treated as an SMT-proposition
and \formula, where \formula is treated as a
type; and this conversion goes \emph{both} ways.
%
For example, let $\formula \equiv p \land (q \lor r)$.
%
Then "flatten" converts from \formula to
$\{\formula\}$ and "expand" the other way,
while this conversion is SMT-aided.
%
\begin{mcode}
    flatten :: p:_ -> q:_ -> r:_-> ({p}, Either {q} {r}) -> {p && (q || r)}
    flatten (pf, Left  qf) = pf &&& qf
    flatten (pf, Right rf) = pf &&& rf

    expand :: p:_ -> q:_ -> r:_ -> {p && (q || r)} -> ({p}, Either {q} {r})
    expand proof | q = (proof, Left  proof)
    expand proof | r = (proof, Right proof)
\end{mcode}

\mypara{Distributing Quantifiers}
% \mypara{Existentials over disjunction}
%
Next, we construct the proof terms needed
to prove two logical properties:
that existentials distribute over disjunctions
% (\ref{eq:ex:or})
and foralls over conjunctions, \ie
% (\ref{eq:ex:and}).
%
\begin{align}
%
  \formula_\exists \ & \equiv (\exists x. p\ x \lor q\ x) \Rightarrow ((\exists x. p\ x) \lor (\exists x. q\ x))
  \label{eq:ex:or} \\
%
  \formula_\forall \ & \equiv (\forall x. p\ x \land q\ x) \Rightarrow ((\forall x. p\ x) \land (\forall x. q\ x))
  \label{eq:ex:and}
%
\end{align}
%
The specification of these properties requires
nesting quantifiers inside connectives and
vice versa.
%
The proof of $\formula_\exists$ (\ref{eq:ex:or})
proceeds by existential case splitting and introduction:
%
\begin{code}
    exDistOr :: p:_ -> q:_ -> (x::a, Either {p x} {q x})
                           -> Either (x::a, {p x}) (x::a, {q x})
    exDistOr _ _ (x, Left  px) = Left  (x, px)
    exDistOr _ _ (x, Right qx) = Right (x, qx)
\end{code}
%
% \PW{Phil suggests to show the proof derivation tree here too}
%
% \mypara{Foralls over conjunction}
% Similarly, we prove that foralls distribute over conjunction:
% %
% $$\formula \equiv \forall x. (p\ x \land q\ x)) \Rightarrow ((\forall x. p\ x)\land (\forall x. q\ x))$$
% %
% The specification of the conclusion now
% requires nesting universal quantification
% over conjunctions.
%
Dually, we prove $\formula_\forall$ (\ref{eq:ex:and}) via
a $\lambda$-abstraction and case spitting
inside the conjunction pair:
%
\begin{code}
    allDistAnd :: p:_ -> q:_ -> (x:a->({p x}, {q x}))
                             -> ((x:a->{p x}), (x:a->{q x}))
    allDistAnd _ _ andx = ( (\x -> case andx x of (px, _) -> px)
                          , (\x -> case andx x of (_, qx) -> qx) )
\end{code}
%
The above proof term exactly corresponds to its natural deduction proof derivation
but using SMT-aided verification can get simplified to the following
%
\begin{code}
    allDistAnd _ _ andx = (pf, pf)
      where pf x = case andx x of (px, py) -> px &&& py
\end{code}

%%\subsubsection{Forall - exists over implication}
%%We prove distribution of foralls over conjunction:
%%$$(\forall x. \exists y. (p\ x \Rightarrow q\ x\ y)) \Rightarrow (\forall x. (p\ x \Rightarrow (\exists y. q\ x\ y))))$$
%%The proof proceeds by forall elimination and existential introduction:
%%
%%\begin{code}
%%forallExistsImpl :: p:(a -> Bool) -> q:(a -> a -> Bool)
%%  -> (x:a -> (y::a, {p x} -> {q x y} ))
%%  -> (x:a -> ({p x} -> (y::a, {q x y})))
%%forallExistsImpl p q f x px
%%  = case f x of
%%      (y, pxToqxy) -> (y, pxToqxy px)
%%\end{code}

% \subsubsection{Properties on user specified domains}

\mypara{Properties of User Defined Datatypes}
%
As \formula can describe properties
of data types like lists, we can prove
properties of such types, \eg that for
every list "xs", if there exists a list
"ys" such that "xs == ys ++ ys" ,then
"xs" has even length.
%
$$\formula \equiv \forall xs. ((\exists ys.\ xs = ys \ \texttt{++}\ ys) \Rightarrow (\exists n. \texttt{len}\ xs = n + n))$$
%
The proof ("evenLen") proceeds by existential elimination and introduction
and uses the "lenAppend" lemma, which in turn uses induction on the
input list and \pbesym to automate equational reasoning.
%
\begin{code}
    evenLen :: xs:[a]->(ys::[a],{xs = ys ++ ys})->(n::Int,{len xs = n+n})
    evenLen xs (ys,pf) = (len ys, lenAppend ys ys &&& pf)

    lenAppend :: xs:_ -> ys:_ -> {len (xs ++ ys) = len xs + len ys}
    lenAppend []     _  = ()
    lenAppend (x:xs) ys = lenAppend xs ys
\end{code}

\mypara{Induction on Natural Numbers}
%
Finally, we specify and verify \emph{induction} on natural numbers:
%
$$
\formula_{ind} \equiv  (p\ 0 \land (\forall n. p\ (n-1) \Rightarrow p\ n) \Rightarrow \forall n. p\ n)
$$
%
The proof proceeds by induction (\eg case splitting).
%
In the base case, "n == 0", the proof calls the left conjunct, which contains a proof of the base case.
%
Otherwise, "0 < n", and the proof \emph{applies} the induction hypothesis to the right conjunct instantiated at "n-1".
%
\begin{code}
    ind :: p:_ -> ({p 0},(n:Nat -> {p (n-1)} -> {p n})) -> n:Nat -> {p n}
    ind p (p0, pn) 0 = p0
    ind p (p0, pn) n = pn n (ind p (p0, pn) (n-1))
\end{code}

\newcommand\fst[1]{\ensuremath{\texttt{fst}\ #1}\xspace}
\newcommand\snd[1]{\ensuremath{\texttt{snd}\ #1}\xspace}

\begin{figure}[t]
%
\setlength\tabcolsep{0pt} % default value: 6pt
\begin{tabular}{ccc}
\multirow{2}{*}{
$
\inference{
    \ch{\env}{\fst{e}}{\formula_1} \\
    \ch{\env}{\snd{e}}{\formula_2}
}{
	\ch{\env}{e}{\tyand{\formula_1}{\formula_2}}
}\land\small{\textsc{\!-I}}$}
& \quad &
$\inference{
	\ch{\env}{e}{\tyand{\formula_1}{\formula_2}}
}{
	\ch{\env}{\fst{e}}{\formula_1}
}
\land\small{\textsc{\!-L-E}}$\\
& \quad & $\inference{
	\ch{\env}{e}{\tyand{\formula_1}{\formula_2}}
}{
	\ch{\env}{\snd{e}}{\formula_2}
}
\land\small{\textsc{\!-R-E}}$
\\[0.12in]
$\inference{
    \ch{\env}{e_1}{\formula_1}
}{
	\ch{\env}{\exprorleft{e_1}}{\tyor{\formula_1}{\formula_2}}
}\lor\small{\textsc{\!-L-I}}$
& \quad &
\multirow{2}{*}{$
\inference{
	\ch{\env}{e}{\tyor{\formula_1}{\formula_2}} \\
	\ch{\env,\bind{x_1}{\formula_1}}{e_1}{\formula} &&
	\ch{\env,\bind{x_2}{\formula_2}}{e_2}{\formula}
}{
	\arraycolsep=0.5pt
	\begin{array}{lll}
	\env \vdash \texttt{case}\ e\ \text{of}\ \{&
		\exprorleft{x_1}\rightarrow e_1; & \\
	& \exprorright{x_2}\rightarrow e_2 \}\ & : \formula	\\
	\end{array}
}\lor\small{\textsc{\!-E}}$}\\
\\
%% Implication
$\inference{
    \ch{\env}{e_1}{\formula_2}
}{
	\ch{\env}{\exprorright{e_2}}{\tyor{\formula_1}{\formula_2}}
}\lor\small{\textsc{\!-R-I}}$ & \quad&
\\[0.12in]
$\inference{
	\ch{\env, \bind{x}{\formula_x}}{e}{\formula}
}{
	\ch{\env}{\lambda x. e}{\formula_x \rightarrow \formula}
}\impIntro$ &
\quad &
$\inference{
	\ch{\env}{e}{\formula_x \rightarrow \formula} &&
	\ch{\env}{e_x}{\formula_x}
}{
	\ch{\env}{e\  e_x}{\formula}
}
\Rightarrow\small{\textsc{\!-E}}
$
\\[0.12in]
%% Forall
$
\inference{
	\ch{\env, \bind{x}{\typ}}{e}{\formula}
}{
	\ch{\env}{\lambda x. e}{(x:\typ \rightarrow \formula)}
}\forallIntro$ &
\quad
&
$
\inference{
	\ch{\env}{e_x}{\typ} &&
	\ch{\env}{e}{(x:\typ \rightarrow \formula)}
}{
	\ch{\env}{e\ e_x}{\formula\subst{x}{e_x}}
}
\forallElim
$
\\[0.12in]
%% Exists
$
\inference{
	\ch{\env}{\fst{e}}{\typ}  &&
	\ch{\env, \bind{x}{\typ}}{\snd{e}}{\formula}  &&
}{
	\ch{\env}{e}{\tyex{x}{\typ}{\formula\subst{x}{\fst{e}}}}
}\exists\small{\textsc{\!-I}}$ &\quad \quad &
$
\inference{
	\ch{\env}{e}{\tyex{x}{\typ}{\formula_x}} &&
	\ch{\env, \bind{x}{\typ}, \bind{y}{\formula_x}}{e'}{\formula}
}{
	\ch{\env}{\texttt{case}\ e\ \texttt{of}\ \{ (x, y) \rightarrow e'\}}{\formula}
}
\existsElim
$
\end{tabular}
%
%% Negation
%%$$
%%\inference{
%%	\ch{\env}{e_x}{\formula} &&
%%	\ch{\env}{e}{\formula \rightarrow \proofterm{\efalse}}
%%}{
%%	\ch{\env}{e\ e_x}{\proofterm{\efalse}}
%%}
%%\neg\small{\textsc{\!-E}}
%%\quad
%%\inference{
%%	\ch{\env, \bind{x}{\formula}}{e}{\proofterm{\efalse}}
%%}{
%%	\ch{\env}{\lambda x. e}{\formula \rightarrow \proofterm{\efalse}}
%%}
%%\neg\small{\textsc{\!-I}}
%%$$
%
\caption{Natural deduction rules for refinement types.
With $[\texttt{fst}|\texttt{snd}]\  e \equiv \texttt{case}\ e\ \texttt{of}\ \{(x_1, x_2)
\rightarrow [x_1 | x_2 ] \}$ .}
\label{fig:ch}
\figrule
\end{figure}


\subsection{Consequences}
%
To summarize, we use the propositions-as-types
principle to make two important contributions.
%
First, we show that natural deduction reasoning
can smoothly co-exist with SMT-based verification
to automate the decidable, quantifier-free portions
of the proof.
%



\begin{comment}

\begin{code}
    andRefine :: b1:Bool -> b2:Bool -> ({b1}, {b2}) -> {b1 && b2}
    andRefine _ _ (b1, b2) = ()
\end{code}
%
Using "andRefine" the user does not have to explicitly eliminate "PAnd".
For instance the "forallAndDistr" proof from~\S~\ref{sec:expressiveness:examples}
simplifies to the following.
%
\begin{code}
    forallAndDistr _ _ andx = (fAndG, fAndG)
      where fAndG = andRefine (f x) (g x) (andx x)
\end{code}

On the other direction, we reify SMT conjunction by using
$\formula_1 \land \formula_2$
as a proof for each conjunct $\formula_1$ and $\formula_2$
that can later be used in natural deduction derivations.
\begin{code}
    andReify :: b1:Bool -> b2:Bool -> {b1 && b2} -> ({b1}, {b2})
    andReify _ _ b = (b, b)
\end{code}
\end{comment}

Second, we show for the first time how
natural deduction proofs can be encoded
in refinement type systems like \toolname
and we expect this encoding to extend,
in a straightforward manner to other
SMT-based deductive verifiers
(\eg \dafny and \fstar).
%
This encoding shows that refinement type systems
are expressive enough to encode any intuitionistic natural deduction proof,
gives a guideline for encoding proofs
with nested quantifiers, and provides
a pleasant implementation of
natural deduction that
% can be used
is pedagogically useful.
